\section{Introdução}

A atividade econômica está associada a externalidades ambientais, entre elas as emissões de dióxido de carbono (CO\textsubscript{2}), que contribuem para as mudanças climáticas. A redução das emissões integra os esforços globais de mitigação das mudanças climáticas \parencite{ipcc2023spm}. A alocação dos esforços de mitigação depende de quais agentes são responsabilizados por quais emissões. 

No âmbito do Protocolo de Kyoto, a contabilização das emissões segue a ótica da produção, onde as emissões são contabilizadas na atividade em que ocorrem \parencite{unfccc2008kyotoManual}. Existem outras óticas de atribuição setorial das emissões de CO\textsubscript{2}. Pela ótica da demanda, as emissões são atribuidas às atividades responsáveis pela demanda final, considerando o encadeamento das emissões diretas e indiretas \parencite{marques2012,millerBlair2009}. As emissões da geração de eletricidade, por exemplo, podem ser registradas no setor elétrico ou atribuídas aos produtos finais cuja produção requer essa eletricidade. A mudança de critério preserva o total contabilizado, mas altera sua distribuição entre as atividades.

A escolha de quem responsabilizar pelas emissões envolve uma dimensão normativa. A ótica da produção identifica a atividade em que a emissão ocorre, mas não determina, por si só, como distribuir os esforços de mitigação. Questões de justiça exigem reflexão e deliberação pública sobre os valores e os fins que devem orientar a vida em sociedade \parencite{sandel2015justica}. Para subsidiar essa deliberação, é necessário compreender como as diferentes óticas de atribuição distribuem as emissões entre as atividades produtivas.

A literatura sobre emissões brasileiras oferece elementos para comparar atribuições entre setores, regiões e cadeias produtivas \parencite{imori2015,sanguinetAzzoni2024,montoyaEtAl2026}. Os totais e os saldos setoriais, entretanto, não mostram, por si sós, quais produtos finais estão associados às emissões de cada atividade nem quais atividades emissoras compõem a atribuição de cada produto final. Também não mostram se os volumes atribuídos se concentram em poucas relações ou se distribuem entre muitas. Examinar essas relações intersetoriais permite complementar a comparação das contas pela produção e pela demanda.

Este trabalho pergunta: \textbf{Como as emissões de CO\textsubscript{2} se distribuem nas relações entre atividades emissoras e setores de produtos finais no Brasil em 2015?} O objetivo é caracterizar essas relações e discutir suas implicações para a coordenação dos esforços de mitigação.

A análise utiliza a extensão ambiental do modelo insumo-produto \parencite{millerBlair2009}, combinando a Matriz de Insumo-Produto (MIP) brasileira de 2015, com 67 atividades \parencite{ibge2015mip}, e coeficientes setoriais de emissão de \textcite{sanguinetAzzoni2024}. As atribuições são representadas em uma rede dirigida e ponderada, adaptando a abordagem de \textcite{huang2019} para relacionar as atividades emissoras aos setores de produtos finais.

O recorte abrange emissões da produção nacional associadas à demanda doméstica e às exportações, excluindo emissões estrangeiras.

Ao relacionar as emissões da produção nacional aos setores de produtos finais, a análise mostra que 39,5\% das emissões atribuídas ao setor da construção são geradas em outras atividades, incluindo as produtoras de materiais e de energia. Esse resultado reforça estratégias de mitigação que combinam a redução das emissões nos setores fornecedores, a eficiência no uso de materiais e critérios de menor emissão nos projetos e nas aquisições da construção \parencite{rissman2020}.
